\chapter{รากฐานฟิสิกส์เกษตรดิจิทัลและสถาปัตยกรรมระบบ}
\label{chap:ch01}

\section*{แผนบริหารการสอนประจำบทที่ 1}
\addcontentsline{toc}{section}{แผนบริหารการสอนประจำบทที่ 1}

\begin{tcolorbox}[
    enhanced, breakable,
    colback=white, colframe=rbruNavy,
    fonttitle=\bfseries\color{white},
    title={1. หัวข้อเนื้อหาประจำบท},
    arc=1.5mm, boxrule=0.8pt, leftrule=3.5mm
]
\begin{enumerate}
    \item ความสำคัญของสมบัติทางกายภาพและเคมีของดินต่อสรีรวิทยาของพืชเศรษฐกิจ
    \item พารามิเตอร์คุณภาพดิน 8 มิติสำหรับการเกษตรแม่นยำ
    \item ทฤษฎีสภาพยอมทางไฟฟ้าและแบบจำลองการแพร่ของไอออนในสารละลายดิน
    \item ข้อจำกัดเชิงโครงสร้างของวิธีการตรวจวัดแบบดั้งเดิมในสนามจริง
    \item สถาปัตยกรรมการประมวลผลสองขั้นตอนแบบ Edge Two-Stage Decoupling
\end{enumerate}
\end{tcolorbox}

\begin{tcolorbox}[
    enhanced, breakable,
    colback=white, colframe=rbruNavy,
    fonttitle=\bfseries\color{white},
    title={2. วัตถุประสงค์เชิงพฤติกรรม},
    arc=1.5mm, boxrule=0.8pt, leftrule=3.5mm
]
เมื่อศึกษาบทเรียนนี้จบแล้ว ผู้เรียนสามารถ
\begin{enumerate}
    \item อธิบายกลไกทางสรีรวิทยาของรากพืชที่ตอบสนองต่อพารามิเตอร์ดินทั้ง 8 ตัวได้อย่างถูกต้อง
    \item อธิบายสมการเชิงประจักษ์ของท็อปป์ (Topp's Equation) และทฤษฎีสภาพยอมทางไฟฟ้าได้อย่างแม่นยำ
    \item วิเคราะห์สาเหตุของความคลาดเคลื่อนจากการเลื่อนลอยตามอุณหภูมิของหัววัดดินได้อย่างเป็นระบบ
    \item ออกแบบสถาปัตยกรรมระบบตรวจวัดดินอัจฉริยะที่ผสานการประมวลผลบนอุปกรณ์สมาร์ตโฟนได้อย่างสมบูรณ์
\end{enumerate}
\end{tcolorbox}

\begin{tcolorbox}[
    enhanced, breakable,
    colback=white, colframe=rbruNavy,
    fonttitle=\bfseries\color{white},
    title={3. วิธีสอนและกิจกรรมการเรียนการสอน},
    arc=1.5mm, boxrule=0.8pt, leftrule=3.5mm
]
\begin{enumerate}
    \item การบรรยายเชิงปฏิสัมพันธ์ร่วมกับการสาธิตปรากฏการณ์ทางฟิสิกส์ในดิน
    \item การอภิปรายกลุ่มย่อยเกี่ยวกับกรณีศึกษาปัญหาดินกรดและดินเค็มในแปลงทุเรียนจังหวัดจันทบุรี
    \item การฝึกปฏิบัติการคำนวณการแปลงสภาพยอมทางไฟฟ้าเป็นความชื้นดินเชิงปริมาตร
\end{enumerate}
\end{tcolorbox}

\begin{tcolorbox}[
    enhanced, breakable,
    colback=white, colframe=rbruNavy,
    fonttitle=\bfseries\color{white},
    title={4. สื่อการเรียนการสอน},
    arc=1.5mm, boxrule=0.8pt, leftrule=3.5mm
]
\begin{enumerate}
    \item สไลด์บรรยายอิเล็กทรอนิกส์และเอกสารคำสอนวิชาฟิสิกส์เกษตรดิจิทัล
    \item ชุดหัววัดดิน 8-in-1 อุตสาหกรรม และตัวอย่างหน้าตัดชั้นดินจำลอง
    \item แอปพลิเคชันจำลองการคำนวณบนสมาร์ตโฟน JC AI Detector
\end{enumerate}
\end{tcolorbox}

\begin{tcolorbox}[
    enhanced, breakable,
    colback=white, colframe=rbruNavy,
    fonttitle=\bfseries\color{white},
    title={5. การวัดและประเมินผล},
    arc=1.5mm, boxrule=0.8pt, leftrule=3.5mm
]
\begin{enumerate}
    \item การประเมินความรู้ผ่านแบบฝึกหัดและคำถามทบทวนท้ายบท
    \item การประเมินทักษะการวิเคราะห์โจทย์ฟิสิกส์เกษตรจากการนำเสนอของกลุ่มผู้เรียน
    \item การสังเกตพฤติกรรมการมีส่วนร่วมและการซักถามในชั้นเรียน
\end{enumerate}
\end{tcolorbox}

\clearpage

% ==============================================================================
% ภาพประกอบเกริ่นนำประจำบทที่ 1 (Introductory Infographic)
% ==============================================================================
\begin{figure}[H]
\centering
\begin{tikzpicture}[
    box/.style={rectangle, rounded corners=2.5mm, draw=rbruNavy, line width=1.0pt, fill=white, drop shadow={shadow xshift=1pt, shadow yshift=-1.5pt, opacity=0.15}, align=center, inner sep=6pt},
    header/.style={rectangle, rounded corners=1.5mm, draw=rbruNavy, fill=rbruNavy, text=white, font=\bfseries\small, align=center, inner sep=4pt},
    arrow/.style={-{Stealth[scale=1.1]}, line width=1.2pt, draw=rbruBlue}
]
    % คอลัมน์ที่ 1: ปัจจัยดินกายภาพและเคมี
    \node[box, fill=cyan!8, minimum width=3.6cm] (col1) {
        \textbf{8 พารามิเตอร์ดิน}\\[3pt]
        \footnotesize
        $\bullet$ อุณหภูมิ \& ความชื้น\\
        $\bullet$ สภาพนำไฟฟ้า (EC)\\
        $\bullet$ ความเป็นกรด-ด่าง (pH)\\
        $\bullet$ ธาตุอาหาร N, P, K\\
        $\bullet$ ดัชนีความสมบูรณ์
    };
    \node[header, above=2pt of col1] {1. แหล่งกำเนิดสัญญาณ};

    % คอลัมน์ที่ 2: ทฤษฎีฟิสิกส์หลักการแรก
    \node[box, fill=rbruGold!10, minimum width=4.0cm, right=1.2cm of col1] (col2) {
        \textbf{สมการหลักการแรก}\\[3pt]
        \footnotesize
        $\bullet$ Topp's Permittivity ($\varepsilon_r$)\\
        $\bullet$ Arrhenius Temp Drift\\
        $\bullet$ Nernst Ionic Potential\\
        $\bullet$ Rhoades EC Conduction
    };
    \node[header, fill=rbruGold!90!black, above=2pt of col2] {2. กลไกฟิสิกส์เชิงลึก};

    % คอลัมน์ที่ 3: สถาปัตยกรรม Edge AI บนสมาร์ตโฟน
    \node[box, fill=green!8, minimum width=3.8cm, right=1.2cm of col2] (col3) {
        \textbf{ระบบ JC AI Detector}\\[3pt]
        \footnotesize
        $\bullet$ USB-C OTG Interface\\
        $\bullet$ RS485 Modbus RTU\\
        $\bullet$ Two-Stage PINN\\
        $\bullet$ Offline Real-Time Dashboard
    };
    \node[header, fill=green!60!black, above=2pt of col3] {3. ปัญญาประดิษฐ์ฝังตัว};

    % ลูกศรเชื่อมโยงกระบวนการ
    \draw[arrow] (col1) -- node[above, font=\scriptsize\bfseries\color{rbruNavy}] {ตรวจวัด} (col2);
    \draw[arrow] (col2) -- node[above, font=\scriptsize\bfseries\color{rbruNavy}] {ชดเชยค่า} (col3);

    % กรอบล่างแสดงเป้าหมายปลายทาง
    \node[rectangle, rounded corners=2mm, draw=rbruNavy!60, fill=rbruNavy!5, line width=0.8pt, minimum width=13.5cm, minimum height=0.9cm, below=0.7cm of col2, align=center] (goal) {
        \small\textbf{เป้าหมายทางวิชาการ:} ยกระดับการประเมินสุขภาพดินแปลงทุเรียนอย่างแม่นยำ ปราศจากของเสียสารเคมี และรู้ผลทันทีในระดับเสี้ยววินาที
    };
    \draw[arrow, dashed] (col2.south) -- (goal.north);
\end{tikzpicture}
\caption{ผังมโนทัศน์เกริ่นนำประจำบทที่ 1: การบูรณาการฟิสิกส์เกษตรดิจิทัล ทฤษฎีหลักการแรก และปัญญาประดิษฐ์บนสมาร์ตโฟน}
\label{fig:ch01_concept_overview}
\end{figure}

\section{ความสำคัญของสมบัติทางกายภาพและเคมีของดิน}

การบริหารจัดการดินในงานเกษตรกรรมแม่นยำ (Precision Agriculture) โดยเฉพาะในแปลงปลูกพืชเศรษฐกิจมูลค่าสูง เช่น ทุเรียนหมอนทอง (\textit{Durio zibethinus} Murr.) ในเขตจังหวัดจันทบุรีและตราด มีความจำเป็นต้องอาศัยข้อมูลสภาพดินที่รวดเร็ว ถูกต้อง และครอบคลุมพารามิเตอร์สำคัญ 8 ประการ

\subsection{สรีรวิทยาของรากพืชและการตอบสนองต่อพารามิเตอร์ดิน}
ระบบรากของทุเรียนเป็นรากที่ต้องการออกซิเจนสูงและมีความไวต่อสภาพความเป็นกรดและความเค็มในดินเป็นอย่างยิ่ง หากดินมีค่าความเป็นกรด-ด่างต่ำกว่า 5.5 ฟอสฟอรัสจะถูกตรึงด้วยเหล็กและอะลูมิเนียม ทำให้พืชไม่สามารถดูดซึมไปใช้ในการพัฒนาตาดอกได้ ในทางกลับกัน หากค่าสภาพนำไฟฟ้า (EC) สูงเกินกว่า 1,500 ไมโครซีเมนต์ต่อเซนติเมตร แรงดันออสโมซิสในสารละลายดินจะสูงกว่าแรงดันภายในเซลล์ราก ส่งผลให้เกิดภาวะขาดน้ำทางสรีรวิทยา (Physiological Drought) และเกิดอาการใบไหม้ร่วงหล่น

\section{พารามิเตอร์คุณภาพดิน 8 มิติสำหรับการเกษตรแม่นยำ}

ระบบตรวจวัดดินอัจฉริยะที่พัฒนาขึ้นมุ่งเน้นการตรวจวัดตัวแปรสำคัญ 8 มิติ ได้แก่
\begin{enumerate}
    \item \textbf{อุณหภูมิดิน (Soil Temperature)} ส่งผลโดยตรงต่อกิจกรรมของจุลินทรีย์และการดูดซึมน้ำของราก
    \item \textbf{ความชื้นในดิน (Soil Moisture Content)} ตรวจวัดสัดส่วนปริมาตรน้ำในดิน (Volumetric Water Content หรือ VWC)
    \item \textbf{สภาพนำไฟฟ้าของดิน (Electrical Conductivity หรือ EC)} บ่งชี้ความเข้มข้นของไอออนเกลือแร่ที่ละลายในดิน
    \item \textbf{ความเป็นกรด-ด่าง (Soil pH)} ควบคุมความสามารถในการละลายของธาตุอาหารหลักและธาตุอาหารรอง
    \item \textbf{ไนโตรเจนที่เป็นประโยชน์ (Available Nitrogen หรือ N)} ส่งเสริมการเจริญเติบโตของกิ่งก้านและใบ
    \item \textbf{ฟอสฟอรัสที่เป็นประโยชน์ (Available Phosphorus หรือ P)} ส่งเสริมการแตกรากฝอยและการสะสมอาหารเพื่อการออกดอก
    \item \textbf{โพแทสเซียมที่เป็นประโยชน์ (Available Potassium หรือ K)} ควบคุมการเปิดปิดปากใบและคุณภาพเนื้อผลไม้
    \item \textbf{ดัชนีความอุดมสมบูรณ์รวม (Soil Fertility Index)} ตัวชี้วัดระดับความสมบูรณ์ของธาตุอาหารโดยรวม
\end{enumerate}

\section{ทฤษฎีฟิสิกส์ของสภาพยอมทางไฟฟ้าและการแพร่ของไอออน}

\subsection{สมการเชิงประจักษ์ของท็อปป์ (Topp's Empirical Equation)}
ดินเป็นตัวกลางรูพรุนหลายวัฏภาคประกอบด้วยอนุภาคดิน ของเหลว (น้ำ) และก๊าซ (อากาศ) ค่าสภาพยอมทางไฟฟ้าสัมพัทธ์ของอากาศเท่ากับ 1 ของอนุภาคดินแข็งเท่ากับ 3 ถึง 5 ขณะที่โมเลกุลน้ำมีค่าสูงถึง 78 ถึง 80 ที่อุณหภูมิห้อง ดังนั้น การวัดสภาพยอมทางไฟฟ้ารวม ($\varepsilon_b$) จึงสะท้อนปริมาณน้ำในดินได้อย่างชัดเจน ตามสมการพหุนามของท็อปป์

\begin{equation}
\theta_v = -0.053 + 0.0292\varepsilon_b - 5.5 \times 10^{-4}\varepsilon_b^2 + 4.3 \times 10^{-6}\varepsilon_b^3
\end{equation}

โดย $\theta_v$ คือ ปริมาณความชื้นเชิงปริมาตร (Volumetric Water Content)

\subsection{สภาพนำไฟฟ้าปรากฏและแบบจำลองของโรดส์ (Rhoades Model)}
ความสัมพันธ์ระหว่างสภาพนำไฟฟ้าปรากฏรวมของดิน ($EC_a$) ที่หัววัดอ่านได้ กับสภาพนำไฟฟ้าของสารละลายในช่องว่างรูพรุน ($EC_w$) อธิบายได้ด้วยแบบจำลองการนำไฟฟ้าขนานกันของโรดส์

\begin{equation}
EC_a = \theta_v \cdot T(\theta_v) \cdot EC_w + EC_s
\end{equation}

โดยที่ $T(\theta_v)$ คือ ฟังก์ชันความคดเคี้ยวของรูพรุน (Tortuosity Factor) และ $EC_s$ คือ สภาพนำไฟฟ้าที่ผิวอนุภาคดิน

\section{ข้อจำกัดเชิงโครงสร้างของวิธีการตรวจวัดแบบดั้งเดิม}

ในสภาพแปลงเกษตรกรรมจริง วิธีการวิเคราะห์ดินแบบดั้งเดิมเผชิญปัญหาเชิงโครงสร้าง 3 ประการสำคัญ ดังแสดงในตารางที่ \ref{tab:comparison_methods}

\begin{table}[H]
\centering
\caption{การเปรียบเทียบสมรรถนะระหว่างวิธีการตรวจวัดดินแบบดั้งเดิมกับระบบ JC Digital Soil AI}
\label{tab:comparison_methods}
\begin{tabularx}{\textwidth}{p{3.0cm}p{3.2cm}p{3.5cm}Y}
    \toprule
    \textbf{คุณลักษณะ} & \textbf{วิธีแล็บเคมีมาตรฐาน} & \textbf{ชุดทดสอบเคมีภาคสนาม} & \textbf{ระบบ JC Digital Soil AI} \\
    \midrule
    ระยะเวลาวิเคราะห์ & 7 ถึง 14 วัน & 15 ถึง 30 นาที & เรียลไทม์ ($< 1.5$ วินาที) \\
    ค่าใช้จ่ายต่อตัวอย่าง & 500 ถึง 1,500 บาท & 50 ถึง 100 บาท & ไม่มีค่าใช้จ่ายซ้ำซ้อน \\
    ความแม่นยำเชิงปริมาณ & สูงมาก ($> 98\%$) & ปานกลางถึงต่ำ ($< 70\%$) & สูง ($> 96\%$ ด้วย AI PINN) \\
    การชดเชยอุณหภูมิ & มีระบบปรับสภาพห้อง & ไม่มี & อัตโนมัติด้วย Deep Learning \\
    ผลกระทบสิ่งแวดล้อม & มีของเสียสารเคมี & เกิดขยะสารเคมีอันตราย & ไร้สารเคมีและเป็นมิตร 100\% \\
    \bottomrule
\end{tabularx}
\end{table}

\section{สถาปัตยกรรมการประมวลผลสองขั้นตอนแบบ Edge Two-Stage Decoupling}

เพื่อแก้ปัญหาความคลาดเคลื่อนสะสมจากการเปลี่ยนแปลงอุณหภูมิดินตามรอบวัน ระบบ \textbf{JC AI Detector} ได้รับการออกแบบสถาปัตยกรรมแบบสองขั้นตอน (Two-Stage Decoupling) ทำงานบนสมาร์ตโฟนโดยตรง

\begin{figure}[H]
\centering
\begin{tikzpicture}[
    node distance=1.6cm,
    block/.style={rectangle, rounded corners=2mm, draw=rbruNavy, fill=rbruNavy!8, thick, minimum width=4.0cm, minimum height=0.9cm, align=center},
    arrow/.style={-{Stealth[scale=1.2]}, thick, draw=rbruBlue}
]
    \node[block, fill=cyan!15] (sensor) {\textbf{หัววัดเซนเซอร์จริง 8-in-1}\\สัญญาณดิบผ่าน USB Type-C OTG};
    \node[block, fill=rbruTeal!20, below=0.8cm of sensor] (stage1) {\textbf{ขั้นตอนที่ 1: First-Principles Physics}\\Topp Permittivity + Arrhenius EC + Nernst};
    \node[block, fill=rbruGold!20, below=0.8cm of stage1] (stage2) {\textbf{ขั้นตอนที่ 2: Edge TinyML PINN}\\Non-linear MLP Error Compensation};
    \node[block, fill=green!20, below=0.8cm of stage2] (ui) {\textbf{JC AI Detector Dashboard}\\แสดงผลสด บันทึก Excel และวิเคราะห์ดิน};

    \draw[arrow] (sensor) -- (stage1);
    \draw[arrow] (stage1) -- (stage2);
    \draw[arrow] (stage2) -- (ui);
\end{tikzpicture}
\caption{ผังการทำงานของสถาปัตยกรรมการประมวลผลสองขั้นตอนแบบ Edge Two-Stage Decoupling}
\label{fig:two_stage_arch}
\end{figure}

\section{สรุปสาระสำคัญประจำบทที่ 1}

การทำความเข้าใจรากฐานฟิสิกส์เกษตรดิจิทัลเป็นบันไดก้าวสำคัญในการออกแบบระบบวิเคราะห์ดินอัจฉริยะ โดยสามารถสังเคราะห์สาระสำคัญได้ดังนี้
\begin{enumerate}
    \item \textbf{บูรณาการ 8 พารามิเตอร์ดิน:} การตรวจวัดดินเพื่อการเกษตรแม่นยำต้องพิจารณาความสัมพันธ์เชื่อมโยงระหว่างสมบัติทางกายภาพ (ความชื้น, อุณหภูมิ, EC) และสมบัติทางเคมี (pH, ไนโตรเจน, ฟอสฟอรัส, โพแทสเซียม) อย่างเป็นองค์รวม เนื่องจากมีผลต่อแรงดันออสโมซิสและการดูดซึมสารอาหารของพืช
    \item \textbf{ข้อจำกัดของหัววัดภาคสนาม:} สัญญาณไฟฟ้าที่ตรวจวัดได้จากหัววัดมักเกิดการเลื่อนลอยตามอุณหภูมิ (Diurnal Temperature Drift) และความไม่เป็นเชิงเส้นอันเนื่องมาจากการเปลี่ยนแปลงสภาพยอมทางไฟฟ้าและการเคลื่อนที่ของไอออนในสารละลายดิน
    \item \textbf{สถาปัตยกรรม Edge Two-Stage Decoupling:} การผสานฟิสิกส์หลักการแรกเข้ากับโครงข่ายประสาทเทียม PINN บนสมาร์ตโฟน ช่วยให้สามารถแยกผลกระทบของอุณหภูมิและความชื้นออกจากค่าความเข้มข้นของธาตุอาหารได้อย่างแม่นยำ โดยไม่ต้องพึ่งพาระบบคลาวด์
\end{enumerate}

ในบทถัดไป (บทที่ 2) จะกล่าวถึงการนำหลักการทางทฤษฎีเหล่านี้ไปสู่การเชื่อมต่อกับฮาร์ดแวร์จริง ผ่านโปรโตคอลอุตสาหกรรม Modbus RTU และพอร์ต USB Type-C OTG บนสมาร์ตโฟน

\section*{คำถามทบทวนท้ายบทที่ 1}
\addcontentsline{toc}{section}{คำถามทบทวนท้ายบทที่ 1}

\begin{enumerate}
    \item จงอธิบายบทบาทของสภาพยอมทางไฟฟ้า ($\varepsilon$) ในการประเมินปริมาณความชื้นในดินตามสมการเชิงประจักษ์ของท็อปป์
    \item เพราะเหตุใดอุณหภูมิดินที่สูงขึ้นในเวลากลางวัน จึงทำให้หัววัดอ่านค่าความชื้นต่ำกว่าความเป็นจริงและค่าสภาพนำไฟฟ้าสูงเกินจริง
    \item จงวิเคราะห์ผลกระทบต่อสรีรวิทยาของรากทุเรียน เมื่อดินมีค่า pH ลดลงต่ำกว่า 5.0 พร้อมระบุกลไกการตรึงธาตุฟอสฟอรัส
    \item จงอธิบายความแตกต่างเชิงสถาปัตยกรรมระหว่างการประมวลผลบนคลาวด์ (Cloud AI) กับการประมวลผลบนอุปกรณ์สมาร์ตโฟน (Edge AI) ในงานเกษตรกรรมแม่นยำ
\end{enumerate}
